% SPDX-License-Identifier: CC-BY-4.0
% Copyright (C) 2026 Fabian Januszewski, Abhijit and the Langlandscape team
% S03-T2: Mazur check.  q = 11: weights (2,3), |Phi^comb| = 6(1/2+1/3) = 5 = num (q-1)/12; T1 table.
% math: mathematics/chapters/s03/1-graphs.tex  prop:s03-phi, lem:s03-elliptic, lem:s03-cyclic, prop:s03-mazur, rem:s03-mazur, ex:s03-mazur
% data: code/results/s03/T1.json (data[q]: h, weights, phi_order, phi_cyclic, eisenstein_all; provenance.params.ells = 2,3,5,7,11,13)
\begin{frame}{Mazur check: $\Phi^{\mathrm{comb}}$ of $\mathcal G_{\mathrm{DR}}$ has order $\operatorname{num}\frac{q-1}{12}$}
  \footnotesize
  \begin{itemize}
    \item \textbf{$q=11$:} $\mathcal X(B_{11})$ has two classes, with weights $w=(2,3)$. By prop.\ s03-phi,
      $|\Phi^{\mathrm{comb}}|=\big(\prod_e w_e\big)\sum_e\tfrac1{w_e}=6\big(\tfrac12+\tfrac13\big)=5=\operatorname{num}\tfrac{11-1}{12}$.
    \item \textbf{General} (prop.\ s03-mazur, \emph{proved modulo cited results}: mass formula, elliptic counts), every prime $q$:
      $\Phi^{\mathrm{comb}}\big(\mathcal G_{\mathrm{DR}}(\widehat{\mathcal O}^\times)\big)\cong \mathbb Z/\operatorname{num}\tfrac{q-1}{12}$, cyclic
      (weights in $\{1,2,3\}$, pairwise coprime; lem.\ s03-elliptic, s03-cyclic).
  \end{itemize}
  \vspace{-1mm}
  \begin{center}\scriptsize
  \begin{tabular}{rrlrcc}\toprule
    $q$ & $h$ & weights $w_e$ & $|\Phi^{\mathrm{comb}}|$ & cyclic & Eisenstein\\\midrule
    $11$  & $2$  & $(2,3)$ & $5$  & yes & yes\\
    $23$  & $3$  & $(1,2,3)$ & $11$ & yes & yes\\
    $37$  & $3$  & $(1,1,1)$ & $3$  & yes & yes\\
    $101$ & $9$  & $(1^8,3)$ & $25$ & yes & yes\\
    $131$ & $12$ & $(1^{10},2,3)$ & $65$ & yes & yes\\\bottomrule
  \end{tabular}
  \end{center}
  \vspace{-1mm}
  \begin{itemize}
    \item \textbf{Mazur comparison} (rem.\ s03-mazur(a), \emph{cited}: Mazur, App.\ Thm.\ A.1(b)): the component group of $J_0(q)$ at $q$ is cyclic of order
      $\operatorname{num}\frac{q-1}{12}$, generated by $(0)-(\infty)$; it agrees with $\Phi^{\mathrm{comb}}$ in order and group structure. The identification of the
      two groups goes through the Deligne--Rapoport dictionary (edges $=$ supersingular points, thickness $=w_e$) and is \emph{not} part of our tree route.
  \end{itemize}
  \vfill
  {\tiny Eisenstein column: $(T_\ell-(1+\ell))$ lands in the image for all $\ell\ne q$, $\ell\le 13$ (computed).
  Sources: \texttt{code/results/s03/T1.json}; prop./rem.\ s03-mazur. Lean: \texttt{Ribet.mazur\_check} formalized modulo \texttt{Hyp.EllipticCount} (and \texttt{Hyp.EichlerMass}).}
\end{frame}
